\section{Regularity in operator norm}
\label{ex-slab-regularity}

To pass from scalar correlation identities to the character construction,
we need regularity in operator norm. Weak continuity of matrix elements
alone would not provide it. The additional ingredient is a local uniform
bound on the operators representing the differentiated insertions.

\begin{lemma}[Operator remainder estimate]\label{ex-operator-remainder}
Let \(U\) be an open subset of a finite-dimensional real normed space
\(E\), let \(H\) be a complex Hilbert space, and let \(\mathcal D\)
be a dense linear subspace. Suppose \(A:U\to B(H)\) and
\(B:U\to\mathcal L_{\mathbb R}(E,B(H))\) satisfy
\[
 D\langle u,A(\cdot)v\rangle(z)[h]=\langle u,B(z)[h]v\rangle
 \quad(u,v\in\mathcal D,\ z\in U,\ h\in E).
\]
If the segment from \(x\) to \(y\) lies in \(U\) and
\(\|B(z)-B_0\|\le C\) along the segment, with \(C\ge0\), then
\begin{equation}\label{ex-operator-remainder-bound}
 \|A(y)-A(x)-B_0[y-x]\|\le C\|y-x\|.
\end{equation}
Consequently, local boundedness of \(B\) implies local Lipschitz
continuity of \(A\); continuity of \(B\) implies \(DA=B\) in
operator norm.
\end{lemma}
\begin{proof}
Fix \(u,v\in\mathcal D\). Along the segment, subtract the linear
function \(z\mapsto\langle u,B_0[z]v\rangle\). The derivative of
the resulting complex-valued function has norm at most
\(C\|u\|\|v\|\). The scalar mean-value inequality yields
\[
 |\langle u,(A(y)-A(x)-B_0[y-x])v\rangle|
 \le C\|u\|\|v\|\|y-x\|.
\]
This inequality also follows directly by composing with the segment and
then with every real-linear functional of norm one on \(\mathbb C\),
using the real mean-value theorem, and taking the supremum. Density
extends the inequality to all \(u,v\in H\). For any bounded operator
\(T\), the supremum of \(|\langle u,Tv\rangle|\) over unit vectors
\(u,v\) is \(\|T\|\): Cauchy--Schwarz gives one inequality and,
when \(Tv\ne0\), the choice \(u=Tv/\|Tv\|\) gives the other.
This proves \eqref{ex-operator-remainder-bound}.

On a sufficiently small ball where \(\|B\|\le C\), use \(B_0=0\)
to obtain the Lipschitz bound. If \(B\) is continuous at \(x\), then
for any \(\epsilon>0\) choose such a ball with
\(\|B(z)-B(x)\|\le\epsilon\). Taking \(B_0=B(x)\) proves that
the remainder is at most \(\epsilon\|y-x\|\), which is precisely
Fr\'echet differentiability with derivative \(B(x)\).
\end{proof}

We apply the lemma to the slab operators of Section~\ref{ex-phases}.
For fixed width \(d>0\), arity \(n\), and labels
\(\boldsymbol v=(v_1,\ldots,v_n)\), write
\(A_{\boldsymbol v}(x)=A_d((x_1,v_1)\cdots(x_n,v_n))\) on
\[
 U_{d,n}=\{x:x_i\ne x_j\ (i\ne j),\quad0<(x_i)_0<d\}.
\]
This is open, since finitely many strict inequalities describe the time
conditions and finitely many diagonals have been excluded.

The two quantitative and differential inputs are as follows. First, for
each grade bound \(D\) and label norm bound \(R\), there are
\(C>0\) and \(K\in\mathbb N\) such that every slab word of length
\(n\), with these label bounds, satisfies
\begin{equation}\label{ex-all-copy-bound}
 \|A_d(w)\|\le C^n(n/r)^{Kn}
\end{equation}
whenever \(0<r\le1\), its points have mutual sup-norm distances at
least \(r\), and each time coordinate lies in \([r,d-r]\).
The constants are independent of \(n,d\) and the positions.
For the empty word the right side is interpreted as one, in agreement
with \(A_d(\varnothing)=E_d\).
Second, for each raw label \(v\) of grade \(d'\) and each coordinate
\(j\), choose a raw label \(\delta_jv\) of grade \(d'+1\) representing
its field derivative. Assume that the smooth scalar correlations obey
the exact slot identity
\begin{equation}\label{ex-raw-slot-identity}
 \partial_{(x_i)_j}S_{\boldsymbol v}(x)
 =S_{\boldsymbol v[i\leftarrow\delta_jv_i]}(x)
\end{equation}
at distinct points, also with any finite list of exterior insertions.
The choice of raw representative need not itself be linear. For a phase
using the original field quotient, the identity includes the assertion
that the chosen representatives have the indicated values in that phase.

\begin{proposition}[Smooth slab operators]\label{ex-smooth-slabs}
Under \eqref{ex-all-copy-bound} and \eqref{ex-raw-slot-identity},
\(A_{\boldsymbol v}:U_{d,n}\to B(\mathcal H)\) is smooth in operator
norm, and
\[
 \partial_{(x_i)_j}A_{\boldsymbol v}(x)
 =A_{\boldsymbol v[i\leftarrow\delta_jv_i]}(x).
\]
\end{proposition}
\begin{proof}
Fix a point \(x\in U_{d,n}\). Choose \(0<r\le1\) below its finitely
many positive time margins and pairwise distances. In a neighborhood of
\(x\), the same margins and distances are at least \(r/2\).
For every fixed label tuple, \eqref{ex-all-copy-bound} therefore gives
a local operator bound. The same holds for every one-slot differentiated
tuple, since there are finitely many labels and their grades and norms
have finite maxima.

Define the real-linear derivative candidate
\[
 B_{\boldsymbol v}(x)[h]
 =\sum_{i=1}^n\sum_{j=0}^3h_{i,j}
    A_{\boldsymbol v[i\leftarrow\delta_jv_i]}(x).
\]
The finite collection of local bounds gives a local bound on
\(\|B_{\boldsymbol v}(x)\|\). For positive words \(p,q\),
\[
 \langle[p],A_{\boldsymbol v}(x)[q]\rangle
 =S((\theta p)\,w(x,\boldsymbol v)\,\tau_dq).
\]
The exterior points lie respectively below zero and above \(d\), so
all configurations remain distinct near \(x\). Differentiating by
\eqref{ex-raw-slot-identity} gives exactly the matrix elements of
\(B_{\boldsymbol v}\). Sesquilinearity extends this statement to
finite word sums, a dense linear subspace. The remainder lemma now
proves norm continuity of \(A_{\boldsymbol v}\), for every tuple.

In particular each differentiated tuple has norm-continuous operators.
Thus the finite sum defining \(B_{\boldsymbol v}\) is norm-continuous.
A second application of the remainder lemma proves
\(DA_{\boldsymbol v}=B_{\boldsymbol v}\). Inductively, if all the
tuple maps are \(C^k\), then each \(B_{\boldsymbol v}\) is \(C^k\)
as a finite sum of fixed coordinate functionals times such maps. The
derivative formula makes all tuple maps \(C^{k+1}\). Starting from
continuity proves smoothness of every finite order, with the stated
coordinate derivatives.
\end{proof}

\begin{corollary}[Norm separability of the compression algebra]
\label{ex-norm-separable}
If the scalar correlations are multilinear in the finite-dimensional raw
label spaces, the hypotheses of Proposition~\ref{ex-smooth-slabs} imply
that \(\mathcal A\) is separable in operator norm.
\end{corollary}
\begin{proof}
Fix a finite basis in every raw label space. At fixed positions, the slab
operator is multilinear in the labels: pairing the purported linear
identity with any two word vectors reduces it to scalar multilinearity;
density then proves equality of bounded operators. Every compressed word
operator is therefore a finite linear combination of ones with basis
labels.

Place a word in positive time and choose a positive rational width
greater than all its time coordinates. Approximate its distinct positions
by rational positions inside this slab. Operator-norm continuity from
Proposition~\ref{ex-smooth-slabs}, followed by the bounded compression
\(A\mapsto P_0A|_{\mathcal H_0}\), proves norm convergence of the
compressed operators. There are countably many rational widths, rational
configurations, arities and finite basis-label tuples. Their compressed
operators thus generate a norm-dense subalgebra. Star polynomials in this
countable set with coefficients in \(\mathbb Q+i\mathbb Q\) form a
countable set dense in the whole norm-closed algebra: approximate the
finitely many generators and coefficients in any given polynomial, and
then use the definition of the closure. This is the required norm
separability.
\end{proof}

The first analytic input, the estimate uniform in the length of the word,
is stronger than the separated-correlation bound used to construct time
translations. The second input is an exact differential identity. The
argument above supplies operator smoothness and norm separability from
these two inputs; it leaves their derivation from the lattice fields to
the construction chapters.
